\documentclass[11pt]{article}
\usepackage[margin=1in]{geometry}
\usepackage[utf8]{inputenc}
\usepackage[T1]{fontenc}
\usepackage{amsmath,amssymb}
\usepackage[dvipsnames]{xcolor}
\usepackage[round]{natbib}
\usepackage{microtype}
\usepackage{booktabs}
\usepackage[colorlinks=true,linkcolor=NavyBlue,citecolor=ForestGreen,urlcolor=NavyBlue]{hyperref}
\setlength{\emergencystretch}{2em}

\newcommand{\FPB}{\mathrm{FPB}}

\title{Catching the Drift Before the Quiz: Prefix-Only Detection and Placebo-Controlled Just-in-Time Intervention Against Skill-Eroding AI Interaction Patterns\thanks{Pre-registered specification draft assembled from the H1.1 discussion pipeline (2026-08-15). \textbf{No experiments have been run}; every numeric parameter is a pre-registered design target with a stated rationale, and every Result slot is an explicit placeholder. Human-subjects work is gated on the IRB and data-use milestones stated in Section~\ref{sec:experiments}.}}
\author{Anonymous}
\date{August 2026}

\begin{document}
\maketitle

\begin{abstract}
AI coding assistance can accelerate task completion while eroding the comprehension developers need to supervise generated code: in a randomized trial of 52 developers learning an unfamiliar Python library \citep{shen2026skill}, the AI-assisted arm scored 17 points lower on a post-task quiz, and the loss tracked six interaction patterns coded in hindsight from complete sessions. Every existing measurement of this harm arrives after the fact --- nothing can act on it while the session is still running. We answer with a gated design: a calibrated prognostic risk that a session, if uninterrupted, terminates in an erosive pattern, estimated from query-stream prefixes alone --- no model internals, no biometrics --- and evaluated against a measured human prefix-labeling ceiling; a frozen model that must pass a sealed, prevalence-anchored exact-bound authorization gate before any live intervention; and, if authorized, a preregistered four-arm trial --- hand-code, unnotified AI, placebo notification, additive comprehension nudge --- that internally replicates the published gap, isolates nudge content from notification reactivity, and preserves the original rubric under blinded grading. Either branch of the gate produces knowledge the field lacks: how early erosive drift is visible at all, or whether catching it preserves learning.
\end{abstract}

\section{Introduction}

AI assistants are now a default companion for programmers working with unfamiliar libraries, which makes one trade newly consequential: the same delegation that speeds the task can starve the comprehension a developer would need to evaluate, debug, and maintain the resulting code. For at least one controlled setting this is no longer speculation. In \citet{shen2026skill}'s randomized trial, 52 developers with Python experience but no exposure to the Trio async library completed tasks with or without AI assistance; the AI-assisted arm scored 17 points lower on a subsequent concept, code-reading, and debugging quiz (50\% vs.\ 67\%, $d = 0.738$), without significant time savings. The harm was not uniform: hindsight coding of the session recordings identified six interaction patterns, and the outcome tracked the pattern, not AI use per se --- full delegation, progressive reliance, and iterative AI debugging scored below 40\%, while generate-then-ask, hybrid code-plus-explanation queries, and purely conceptual questioning scored 65--86\%.

Yet a central question about this harm remains unexamined: every measurement of it arrives after the fact. The base trial labels complete recordings; observational work detects an ``AI-prompt signature'' only retrospectively and at cohort level \citep{yao2026scaffold}; EEG and recall studies measure downstream states, not the interaction stream \citep{kosmyna2025brain}. Nobody knows how early an erosive session announces itself in behavior --- to a statistical procedure or even to an expert human reading the transcript --- and as long as the drift is only nameable after the quiz, assistants can do nothing about it at the moment it matters, and the people incurring the harm cannot see it at decision time. The question stated by the base paper's framing but attempted by no one is whether the early behavioral signature of an erosive session is recognizable from a partial session prefix, and whether acting on that recognition preserves learning.

Two obstacles make this harder than it looks. First, the pattern taxonomy is a hindsight summary of whole sessions: a learner may be genuinely pattern-indeterminate at query three, so a prefix-time ``diagnosis'' of the eventual pattern is partially ill-posed by construction. The detector must instead estimate a prognostic risk --- the probability that the session, if uninterrupted, terminates in an erosive pattern --- and the honest ceiling on any such estimate is itself unknown, so we measure it directly by asking annotators to label prefixes blind to session endings. Second, an intervention triggered by detection can ``work'' for the wrong reasons: interruption alone changes behavior \citep{bucinca2021trust}, and an intervention that degrades the assistant merely re-derives the hand-coding arm. Worse, intervening on live participants on the strength of a noisy small-sample precision estimate is an ethical failure before it is a statistical one.

Our design turns each obstacle into a specific, pre-registered mechanism. The risk model reads only what a logging proxy sees --- query text, timestamps, responses, paste and error events --- so the instrument is deployable by any party, including auditors who do not own the model. Its output is converted into at most one timed alert per session by a first-crossing trigger whose risk value is frozen before any message is delivered. No live participant is intervened on unless a frozen model and cutoff first pass a sealed authorization gate: an exact Clopper--Pearson bound \citep{clopper1934use} on alert precision must clear the measured erosive base rate by a pre-registered margin, with a minimum alert count and an explicit annoyance budget --- not an arbitrary constant threshold. And the live trial that a gate pass authorizes carries its own validity checks: hand-code and unnotified-AI arms internally replicate the published contrast before any intervention contrast is interpreted, and a timing-, length-, and presentation-matched placebo notification separates what the nudge says from the fact that it interrupts.

The program is deliberately serialized so that either branch of the gate produces knowledge the field lacks. A gate failure still delivers the first quantification of how early erosive patterns are visible in the behavior stream --- to humans and to a leakage-controlled statistical procedure --- which is a publishable characterization with direct design implications. A gate pass additionally delivers the first closed just-in-time loop outside purpose-built tutoring systems whose sensed state is an AI-interaction pattern and whose distal outcome is measured skill formation, on an instrument directly comparable to the published baseline.

In summary, we contribute:
\begin{itemize}
  \item A prognostic, prefix-only risk model for skill-eroding AI interaction, validated with developer-disjoint cross-fitting against prevalence baselines and a measured \textbf{human prefix-labeling ceiling} --- itself a new scientific quantity: how early are these patterns visible at all?
  \item A \textbf{decision-analytic, sealed authorization gate} that converts small-sample alert precision into an auditable, prevalence-anchored go/no-go rule, pre-registered as the ethical precondition for any live intervention.
  \item A gate-conditional, \textbf{four-arm randomized evaluation} that internally replicates the published hand-vs-AI skill-formation gap, estimates the effect of a one-shot additive comprehension nudge, and isolates nudge content from notification reactivity via a matched placebo.
  \item A \textbf{scoped novelty position}: we import, and credit, the just-in-time adaptive intervention vocabulary \citep{nahumshani2018jitai}, the intelligent-tutoring closed-loop tradition \citep{corbett1995knowledge} --- including its gaming-the-system detect-and-intervene line --- and the learning-analytics early-warning detector shape \citep{gardner2018mooc}; the contribution is the transplant of that loop to the setting where the assistant itself is the erosion vector, during authentic work, sensed through zero-hardware query-log features.
\end{itemize}

\section{Problem Definition}
\label{sec:problem}

\paragraph{Formal setup.}
For developer $i$, a session is an ordered query stream $Q_i = (q_{i1}, \ldots, q_{iT_i})$ produced while working with an AI assistant on an unfamiliar-library task, where each $q_{it}$ carries its text, timestamp, the assistant's response, and any paste/error events. Let $L_i \in \{0,1\}$ be the hindsight label --- erosive- versus protective-terminal --- assigned to the \emph{complete} session by the published six-pattern coding protocol \citep{shen2026skill}, collapsed to a binary. For prefix length $k$, write $Q_i^{(k)} = (q_{i1}, \ldots, q_{ik})$.

\paragraph{Stage 1 (offline, detection).}
Learn $f: Q^{(k)} \to [0,1]$ producing a calibrated prognostic risk $r_{ik} = f(Q_i^{(k)}) \approx P(L_i = 1 \mid Q_i^{(k)})$ using only prefix-observable features --- no model internals, no future queries, no full-session statistics. With cutoff $c$, the alert rule is $\hat{L}_{ik} = \mathbb{1}[r_{ik} \geq c]$. Evaluation reports precision and recall at prefix lengths $k \in \{1,3,5,10\}$ and elapsed-time checkpoints $\{10, 20, 30\}$ minutes, against (a) the prevalence baseline and (b) the human prefix ceiling: the agreement of annotators who label the same prefixes blind to session endings.

\paragraph{Authorization gate.}
Let $\hat{p}$ be the erosive-label prevalence on development data. One final $(f, c)$ is frozen and hashed, then scored exactly once, at the half-session checkpoint, on a sealed developer-disjoint gate set. The live study is authorized iff
\begin{equation}
G \;=\; \mathbb{1}\!\left[\; \underline{\pi}_{90} \geq \hat{p} + 0.10 \;\;\wedge\;\; a \geq 5 \;\;\wedge\;\; \FPB \leq 0.3 \;\right] = 1,
\label{eq:gate}
\end{equation}
where $\underline{\pi}_{90}$ is the one-sided 90\% Clopper--Pearson lower bound \citep{clopper1934use} on gate-set alert precision, $a$ the number of alerted sessions, and $\FPB$ the false-positive burden --- expected alerts per protective-terminal session (the annoyance budget). A design \emph{target} precision of 0.70 is retained only if the measured base rate endorses it; the gate itself is anchored to prevalence, never to a constant.

\paragraph{Stage 2 (live, intervention; only if $G = 1$).}
Participants are randomized to four concurrent arms $A_i \in \{H, U, P, N\}$: hand-code ($H$), unnotified AI assistance ($U$), placebo-trigger ($P$), and active-nudge ($N$). The frozen classifier runs silently in all AI arms; in $P$ and $N$, the first crossing $r_{ik} \geq c$ after a floor of $k \geq 3$ triggers, at most once per session, either a matched content-free notification ($P$) or an additive comprehension nudge ($N$); the risk value is frozen before any message is delivered, so flagged subgroups are defined pre-treatment in every AI arm. With $Y_i$ the blinded, original-rubric quiz score and $\mu_a = \mathbb{E}[Y_i \mid A_i = a]$, the pre-registered estimands are:
\begin{itemize}
  \item \textbf{Replication (validity precondition):} $\Delta_{\mathrm{rep}} = \mu_U - \mu_H$, expected negative, directionally replicating the published 17-point gap;
  \item \textbf{Primary:} $\Delta_{\mathrm{content}} = \mu_N - \mu_P$ (nudge content beyond notification reactivity), intention-to-treat and within the pre-trigger flagged subgroup;
  \item \textbf{Secondary:} $\Delta_{\mathrm{react}} = \mu_P - \mu_U$ and $\Delta_{\mathrm{policy}} = \mu_N - \mu_U$;
  \item \textbf{Process endpoints:} detection lag (query index, wall-clock, and fraction of session remaining at first alert) and guardrails --- task completion and time-on-task in $N$ within 10\% of $U$.
\end{itemize}

\paragraph{Why this is hard.}
The label $L_i$ exists only at session end, so training $f$ against it presumes the terminal label was ``already true'' at query $k$ --- false for learners who drift or self-correct; this is label noise by construction, and it bounds achievable prefix precision from above by a ceiling nobody has measured. The observable stream is thin: the taxonomy that defines $L_i$ was built from rich qualitative reading, while $f$ sees only log features, so signal sufficiency is an open empirical question, not an assumption. Any effect of a triggered message on $Y$ confounds content with interruption reactivity unless the trigger policy, timing, and presentation are matched across $P$ and $N$. The authorization decision must be made from tens of sessions, where point estimates of precision are noise --- only exact small-sample bounds, sealed data, and a pre-enumerated decision table make $G$ trustworthy. And the entire comparison inherits a fragile base effect ($n = 52$, one lab), which is why $\Delta_{\mathrm{rep}}$ is carried inside the design rather than assumed from the literature.

\section{Method Overview}

The study is one measurement-then-intervention pipeline with six cooperating parts. Query streams enter as prefixes; a \textbf{prognostic risk model} turns them into calibrated erosive-terminal risk after every query. A dual-purpose \textbf{annotation harness} supplies both the hindsight labels the model trains against and the human prefix ceiling that bounds what any prefix-time estimate can achieve. A \textbf{first-crossing trigger} converts the running risk into at most one timed, logged alert per session, freezing the risk value before any message is sent. A \textbf{sealed authorization gate} decides, from exact small-sample bounds on sealed data, whether the frozen detector is good enough to justify intervening on people. If it is, an \textbf{additive comprehension nudge with a matched placebo} delivers the intervention in a way that makes its mechanism attributable --- content versus reactivity --- and a \textbf{four-arm assignment scheme} keeps the published hand-vs-AI contrast internally re-tested alongside the intervention contrasts. The parts are serialized left to right: no trigger fires with an unfrozen model, no participant is nudged without a gate pass, and no intervention contrast is interpreted before the replication check.

\section{Methodology}

\subsection{Prefix-Only Prognostic Risk Model}

\textbf{Challenge.} Ground truth is a hindsight label on the complete session; a learner can be genuinely pattern-indeterminate early, so prefix-time six-way diagnosis is partially ill-posed, and the informative behaviors are partly semantic while a logging proxy sees only thin features.

\textbf{Insight.} Reframing the target as \emph{terminal-pattern risk} turns early indeterminacy into probability mass rather than label noise --- a calibrated model can honestly output 0.55 --- and a landmark structure over cumulative prefix features yields a per-query estimate that cannot leak future information by construction.

\textbf{Method.} Two feature tiers are computed cumulatively at each prefix: deterministic log features (query count and length, inter-query gaps, elapsed time, code-block presence, code-paste-then-error-paste cycles, follow-up-after-code presence, user-authored diagnostic detail), and LLM-tagged query intents (complete-solution request / explanation-or-testing request / hybrid / conceptual question) produced by a frozen prompt with a 15\% human audit at $\kappa \geq 0.75$. A session-weighted elastic-net logistic landmark model with Platt calibration estimates $r_{ik}$ after each query; sessions contribute equal total weight regardless of prefix count; hyperparameters and calibration are selected by nested leave-one-developer-out cross-fitting; no deep models at this $n$. The cutoff $c$ is selected on fold-train data only, targeting the gate's operating point subject to the annoyance budget. A deterministic-only ablation answers how thin the signal can be rather than assuming it.

\subsection{Hindsight-Label and Human Prefix-Ceiling Harness}

\textbf{Challenge.} Any classifier's prefix-time performance is bounded above by how early the patterns are visible \emph{at all} --- a quantity nobody has measured --- and the labels themselves come from a qualitative protocol whose reliability at this site must be demonstrated, not assumed.

\textbf{Insight.} The same annotation apparatus can produce both: protocol-coded terminal labels (the training target) and blind prefix labels (the ceiling), making the ceiling simultaneously the honest denominator for the model, a standalone scientific result, and the cheapest possible de-risking of the gate --- if humans cannot call the pattern from half a session, no classifier can.

\textbf{Method.} Two annotators code all complete sessions with the published six-pattern protocol; the erosive/protective binary is primary ($\kappa \geq 0.7$ required; disputes adjudicated; six-way classification kept descriptive). Separately, annotators label a 30-session subset of \emph{prefixes} at $k \in \{1,3,5,10\}$, blind to session endings; their agreement with terminal labels, per $k$, is the human prefix ceiling. A prefix-time validity check verifies that flagged-at-$k$ sessions differ in \emph{subsequent} behavior (post-$k$ feature trajectories), so flags are prognostic rather than correlates of prefix quirks.

\subsection{First-Crossing Trigger and Detection-Lag Instrumentation}

\textbf{Challenge.} A single fixed checkpoint can show that signal exists by minute 30, but not how much session remains when it appears --- and ``just-in-time'' is a timing property: an alert with most of the session already spent supports no intervention worth the name.

\textbf{Insight.} Making the trigger event-driven --- the first crossing of a frozen cutoff --- turns detection lag into a measured distribution, and freezing the risk value before any message is delivered is what keeps flagged-subgroup comparisons causal: the subgroup is defined by pre-treatment information in every arm.

\textbf{Method.} $r_{ik}$ is recomputed after every query; the alert is the first crossing of the frozen $c$ after a floor of three queries; at most one alert fires per session (bounding the loop's non-stationarity); the risk value and feature snapshot are logged before delivery. Detection lag is recorded as query index, wall-clock time, and the fraction of the session's eventual query count still ahead.

\subsection{Decision-Analytic Sealed Authorization Gate}

\textbf{Challenge.} The go/no-go decision to intervene on live participants must be made from tens of sessions, where a precision point estimate is noise, a fixed-constant bar is either unpassable by construction or vacuously easy depending on the base rate, and any post-hoc freedom in threshold or cutoff choice invalidates the guarantee.

\textbf{Insight.} At this $n$ only exact binomial bounds are honest --- a go/no-go discipline imported from Simon-style two-stage phase~II clinical-trial designs, here gating a detector's licence to intervene rather than a treatment effect --- and anchoring the bar to the \emph{measured} prevalence --- pass only if the lower confidence bound on alert precision beats the base rate by a margin, with enough alerts to mean something and a bounded false-positive burden --- makes the gate simultaneously passable by a genuinely enriched detector and unpassable by a lucky one. Sealing plus a pre-enumerated pass table removes every degree of post-hoc freedom.

\textbf{Method.} After development analyses close, one final $(f, c)$ is fitted on all development sessions and frozen: serialized and hashed together with the feature specification, exclusion rules, and the fully enumerated (alerts, true-positives) pass table. The sealed, developer-disjoint gate set (at least 20 sessions, untouched during development) is scored exactly once --- a single pass that evaluates the half-session criterion and logs each session's full per-query risk trajectory, from which all pre-registered read-outs (the gate decision and the lead-time distribution) are computed without ever re-scoring sealed data --- and the criterion of Eq.~\eqref{eq:gate} is applied mechanically. Zero alerts is a pre-registered no-go. One re-gate is permitted after a pre-registered narrowing to delegation-pattern-only detection; a second failure retires the real-time premise as scoped and redirects the topic to retrospective session-audit tooling.

\subsection{Additive Comprehension Nudge with Matched Placebo}

\textbf{Challenge.} A triggered intervention can improve quiz scores for the wrong reasons: interruption alone increases effort \citep{bucinca2021trust}, and withholding answers re-derives the hand-coding arm --- either way the mechanism claim is spurious, and an intervention that tanks task completion has ``worked'' vacuously.

\textbf{Insight.} Matching the placebo to the active nudge in trigger policy, timing, length, and presentation leaves exactly one difference --- whether the appended content does cognitive work --- so an effect that survives the placebo contrast and the completion guardrails is attributable to content acting on an ongoing session.

\textbf{Method.} On trigger, the assistant answers exactly as it otherwise would, then appends one targeted comprehension question keyed to the just-delivered code, instantiated by an LLM from a vetted template bank with a deterministic template fallback; nothing is ever withheld. The placebo fires under the identical policy with content-free text. One nudge per session; 100\% of pilot firings are audited for relevance and timing (intervention-integrity metrics). Pre-registered guardrails require task-completion rate and time-on-task in the active arm to remain within 10\% of the unnotified arm.

\subsection{Internally Replicated Four-Arm Design}

\textbf{Challenge.} The headline comparability claim --- results directly interpretable against the published 50\%/67\% baseline --- rests on a base effect measured once, at $n = 52$, in one lab; a design that assumes it invites collapse, and one that cannot re-test it forfeits the comparison.

\textbf{Insight.} Carrying the published trial's own two arms inside the new trial converts the base effect's fragility from an inherited assumption into a measured result, and a pre-registered interpretation tree makes a replication failure informative rather than fatal.

\textbf{Method.} Participants are randomized, stratified by screening-quiz score, to hand-code, unnotified-AI, placebo-trigger, or active-nudge arms; the frozen classifier runs silently in all AI arms; AI-arm allocation and the algorithmic trigger are concealed until debrief (IRB-approved partial disclosure). Quizzes are graded by two arm-blinded graders with the original rubric ($\kappa$ reported). Analysis is ANCOVA with baseline covariates (experience, screening score), multiple imputation \citep[50 datasets, Rubin's rules;][]{rubin1987multiple}, and delta-adjusted MNAR tipping-point sensitivity. The interpretation tree is pre-registered: if the replication contrast fails directionally, all intervention contrasts are reported as estimation-only and the replication failure is itself a headline finding.

\section{Experiments}
\label{sec:experiments}

\subsection{Experimental Setup}

\paragraph{Data.}
\emph{Source A (pursued, not depended on):} the base trial's 52 AI-arm session transcripts under a data-use agreement with the authors; the request is sent in week 1 with a pre-registered answer deadline of 2026-09-15. \emph{Source B (fallback, auto-starts 2026-09-16):} fresh collection of 40--50 AI-assisted sessions on the reconstructed task under a lightweight transcript-collection IRB protocol, extending to ${\sim}60$ sessions if the data-use agreement fails outright --- sized so that the offline split holds in every branch; these sessions double as the internal replication's first data, and every deviation of reconstructed materials from the original task and quiz is documented. All available sessions are pooled and split \textbf{by developer} into a development set ($\geq 40$ sessions) and a sealed gate set ($\geq 20$ sessions) via an outcome-blind blocked hash balancing task family, experience, and source --- correcting an earlier 20-development/32-gate allocation that starved the model and made the gate's precision requirement statistically unreachable. Sessions collected beyond the 40-session development minimum are pre-committed to the gate set: at the 20-session minimum the pass table demands observed precision near 0.90 at 9--10 alerts, and a larger gate set is the only pre-registered lever that relaxes the demanded cell without touching the bar itself. \emph{Live trial:} four arms, $n = 120$ (30/arm), 144 recruited (20\% over-recruitment), Prolific/Upwork, 1+ year Python, no Trio experience, screening quiz with embedded familiarity traps, \$100--150 per 2--2.5\,h session.

\paragraph{Comparators.}
Controls and references --- this is a measurement-and-intervention study; no research question is framed as ``our procedure beats baseline X.'' The fold-train prevalence-only predictor is the sanity reference for every probabilistic claim. The human prefix ceiling is the upper reference for prefix-time performance. The hand-code and unnotified-AI arms reproduce the published trial's own two arms; the placebo-trigger arm isolates notification reactivity.

\paragraph{Metrics.}
Alert precision/recall and enrichment over prevalence at prefix lengths $k \in \{1,3,5,10\}$ and elapsed-time checkpoints; session-weighted Brier score and calibration diagnostics; human-ceiling agreement per $k$; gate quantities (alert count, true positives, Clopper--Pearson bound, false-positive burden); detection lead time; blinded original-rubric quiz scores by arm; task-completion and time-on-task guardrails; nudge-firing integrity.

\paragraph{Ethics, IRB, and timeline.}
Two IRB protocols are pipelined: transcript collection (minimal risk, consented logging; submitted immediately) and the live RCT (mild information asymmetry --- participants know sessions are logged and an AI assistant is involved, not which arm or that nudges are algorithmically triggered until debrief; submitted 2026-11-15 for a Feb 2027 start under expedited review). The model, cutoff, pass table, nudge templates, and analysis plan are frozen, hashed, and pre-registered (OSF) before the first live participant. Gate calendar: G0 (labeled offline set, annotator $\kappa$) 2026-10-31; G1 (sealed gate) 2027-01-15, one re-gate absorbing to 2027-03-01; pilot Feb--May 2027; G2 (replication + guardrails) 2027-05-31; scale-up Jun--Aug 2027; CHI 2028 submission ${\sim}$Sep 2027 with CSCW/AIES fallbacks. A data-free FAccT 2027 position form (deadline 2026-11-03) argues the normative case using only the published evidence base and this gated design. Budget envelope: \$35--50k participants and annotation, 6--8 person-months, negligible compute.

\paragraph{Research questions.}
\begin{enumerate}
  \item \textbf{RQ1} --- How early are erosive patterns humanly visible from prefixes (the ceiling)?
  \item \textbf{RQ2} --- Does prefix-time alert precision clear the prevalence-anchored margin at half-session?
  \item \textbf{RQ3} --- Do calibrated risk scores beat prevalence-only prediction at locked checkpoints?
  \item \textbf{RQ4} --- Do deterministic, zero-hardware features alone retain most of the signal?
  \item \textbf{RQ5} --- Does the frozen detector pass the sealed authorization gate?
  \item \textbf{RQ6} --- Does the first alert fire with enough session remaining to act (lead time)?
  \item \textbf{RQ7} --- Does the published hand-vs-AI gap internally replicate?
  \item \textbf{RQ8} --- Does nudge content improve quiz scores beyond notification reactivity, within guardrails?
\end{enumerate}

\subsection{RQ1: How Early Are Erosive Patterns Humanly Visible? (the Ceiling)}

\textbf{Setup.} Annotators label 30 development-session prefixes at $k \in \{1,3,5,10\}$, blind to session endings; agreement with protocol-coded terminal labels is computed per $k$. Sample size is a characterization target: 120 prefix judgments per annotator, sized for a $\pm 0.15$ CI half-width (at 90\% confidence) on per-$k$ precision at the implied base rate.

\textbf{Expected outcome.} The human ceiling clears prevalence $+\,0.10$ precision at the half-median session prefix but not at $k = 1$: the patterns become visible somewhere between the first query and mid-session.

\textbf{Falsification condition.} Refuted if half-session human precision fails to clear prevalence $+\,0.10$ --- the patterns would not be humanly visible early, bounding the real-time premise away; the ceiling then becomes the headline finding and the classifier questions are reported as characterization only.

\textbf{Hypothesis (plain language).} Erosive-terminal sessions are recognizably different, to expert eyes, by the middle of the session.

\textbf{Metrics.} Human precision/recall vs.\ terminal labels per $k$; inter-annotator $\kappa$ on prefixes.

\textbf{Result.} \emph{[placeholder --- fill in after experiments are run]}

\subsection{RQ2: Does Prefix-Time Alert Precision Clear the Prevalence-Anchored Margin?}

\textbf{Setup.} Development set only ($\geq 40$ sessions, ${\sim}16$--24 erosive under the design assumption checked at G0); nested leave-one-developer-out cross-fitting; cluster bootstrap over developers with the full pipeline re-run per replicate ($B = 2{,}000$); precision/recall at $k \in \{1,3,5,10\}$ and 10/20/30-minute checkpoints.

\textbf{Sample-size rationale.} At fewer than ${\sim}40$ development sessions the bootstrap lower bound cannot clear a $+\,0.10$ margin even for a genuinely enriched detector --- the feasibility failure that sank the prior parameterization --- hence the dual-source data plan.

\textbf{Expected outcome.} The one-sided 90\% cluster-bootstrap lower bound on alert precision exceeds prevalence $+\,0.10$ at some $k$ no greater than half the median session length; not at $k = 1$.

\textbf{Falsification condition.} Refuted if the lower bound fails the margin at every $k \leq$ half-median session --- the detector would add nothing over a base-rate guess, a pre-registered no-go before any sealing.

\textbf{Hypothesis (plain language).} A leakage-controlled statistical procedure can identify erosive-terminal sessions from the first half of the query stream more precisely than the base rate.

\textbf{Metrics.} Precision, recall, alert count, enrichment over prevalence per $k$ and per checkpoint; gap to the RQ1 ceiling.

\textbf{Result.} \emph{[placeholder --- fill in after experiments are run]}

\subsection{RQ3: Do Calibrated Risk Scores Beat Prevalence-Only Prediction?}

\textbf{Setup.} As RQ2; session-weighted Brier scores at locked 10/20/30-minute checkpoints against the fold-train prevalence-only predictor; Holm correction \citep{holm1979simple} across the three checkpoints. Pre-registered ablations: regularization on/off, calibration on/off, checkpoint-by-index vs.\ by-time.

\textbf{Expected outcome.} Relative Brier improvement with a cluster-bootstrap one-sided 90\% lower bound above 0.10 at the 20- and 30-minute checkpoints; the 10-minute checkpoint may be inconclusive.

\textbf{Falsification condition.} Refuted if no checkpoint clears the Holm-adjusted bound --- the risk scores would carry no calibrated information beyond the base rate.

\textbf{Hypothesis (plain language).} The model's probabilities are better forecasts of how a session ends than knowing only how often sessions end badly.

\textbf{Metrics.} Session-weighted Brier per checkpoint; calibration slope and intercept; reliability diagrams.

\textbf{Result.} \emph{[placeholder --- fill in after experiments are run]}

\subsection{RQ4: Can the Features Be Thin? (Zero-Hardware Sufficiency)}

\textbf{Setup.} Re-run the full RQ3 pipeline with deterministic log features only (no LLM intent tags); compare relative Brier improvement to the full two-tier model at the half-session checkpoint; cluster-bootstrap CI on the retained fraction. This question operationalizes the open proxy-validity issue inherited from gaze-based reliance sensing \citep{cao2022reliance}: the validated real-time proxy is hardware-grade, and whether text-only features match it is untested.

\textbf{Expected outcome.} The deterministic-only model retains at least 50\% of the full model's relative Brier improvement --- the signal does not hinge on semantic intent tagging --- while the full model remains better.

\textbf{Falsification condition.} Refuted if the retained fraction's confidence interval falls below 0.5 --- the detector's signal would depend on the LLM tagger, and the ``any logging proxy can deploy this'' claim narrows to proxies with tagging capability (scope note pre-committed).

\textbf{Hypothesis (plain language).} Counting and timing behaviors --- not reading their meaning --- carries most of the early-warning signal.

\textbf{Metrics.} Retained-fraction ratio with CI; per-tier feature importances (descriptive).

\textbf{Result.} \emph{[placeholder --- fill in after experiments are run]}

\subsection{RQ5: Does the Frozen Detector Pass the Sealed Authorization Gate?}

\textbf{Setup.} One final model-and-cutoff pair is frozen and hashed with the feature specification, exclusion rules, and the fully enumerated pass table; the sealed developer-disjoint gate set ($\geq 20$ sessions) is scored exactly once --- a single pass evaluating the half-session criterion and logging per-query risk trajectories --- and the criterion of Eq.~\eqref{eq:gate} is applied mechanically. At $n = 20$ and prevalence ${\approx}\,0.5$, the criterion is passable at, e.g., 8 true positives among 9 alerts (Clopper--Pearson 90\% lower bound ${\approx}\,0.632 > 0.60$) or 9 of 10 (${\approx}\,0.663$) --- reachable by a genuinely enriched detector --- whereas 7 of 8 fails (0.594), and a fixed 0.70 lower-bound rule would be cleared only by near-perfect cells such as 8 of 8 (0.750). The minimum-set cells are demanding --- 8/9 and 9/10 are observed precisions of 0.89--0.90 --- which is why surplus sessions are pre-committed to the gate set: at ${\sim}15$ alerts the demanded cell relaxes to 12 true positives among 15 (90\% lower bound ${\approx}\,0.607$).

\textbf{Expected outcome.} Conditional, not confident: at the pre-registered bar, a detector whose true precision merely clears the RQ2 margin (0.60--0.75) passes a 9--10-alert gate with low probability, so a first-attempt pass is expected only if the detector's true precision lands well above the margin or the surplus-to-gate-set pre-commitment enlarges the alert count; the design's honest expectation is a pass within the two permitted attempts --- lower bound above prevalence $+\,0.10$, at least 5 alerts, false-positive burden within the 0.3 annoyance budget --- with the characterization paper as the pre-committed product otherwise.

\textbf{Falsification condition.} The gate is the falsification instrument: failing it twice (including the one permitted re-gate after narrowing to delegation-only detection) refutes the deployable-detector claim at this scale and scope, and the pre-committed failure product is the RQ1--RQ4 characterization paper.

\textbf{Hypothesis (plain language).} The detector's precision advantage survives contact with sealed data it has never influenced.

\textbf{Metrics.} Alert count, true positives, precision, Clopper--Pearson bound, false-positive burden.

\textbf{Result.} \emph{[placeholder --- fill in after experiments are run]}

\subsection{RQ6: Does the First Alert Fire with Enough Session Remaining to Act?}

\textbf{Setup.} First-crossing times (floor $k \geq 3$) are read from the per-query risk trajectories logged during the gate set's single scoring pass --- sealed data is never re-scored; lead time is the fraction of the session's eventual query count still ahead at the first true alert; one-sided sign test on the median. The 40\% floor is an authored design target: an alert with most of the session spent cannot support a just-in-time intervention, and 40\%-remaining leaves the nudge roughly half the working session to act on.

\textbf{Expected outcome.} Median lead time among true alerts is at least 40\% of the session's eventual query count.

\textbf{Falsification condition.} Refuted if the median true alert fires with less than 40\% of the session remaining --- detection would be real but too late, and the intervention claim would be bounded by timing rather than accuracy; the live trial's detection-lag endpoint re-checks this in vivo.

\textbf{Hypothesis (plain language).} When the alarm rings, most of the fire is still preventable.

\textbf{Metrics.} Lead-time distribution (query-index and wall-clock forms); later, per-arm lag in the live trial.

\textbf{Result.} \emph{[placeholder --- fill in after experiments are run]}

\subsection{RQ7: Does the Published Hand-vs-AI Gap Internally Replicate?}

\textbf{Setup.} Hand-code vs.\ unnotified-AI arms of the live trial; blinded original-rubric grading by two graders ($\kappa$ reported); one-sided ANCOVA contrast at $\alpha = 0.05$ with baseline covariates; multiple imputation with Rubin's rules \citep{rubin1987multiple}; MNAR tipping-point sensitivity. At 30/arm, power is at least ${\sim}80\%$ for the published $d = 0.738$ (two-sample normal approximation ${\approx}\,0.89$) --- the one contrast the pilot is confirmatorily powered for, by design.

\textbf{Expected outcome.} The unnotified-AI arm scores directionally and significantly lower than the hand-code arm.

\textbf{Falsification condition.} Refuted if the unnotified-AI arm scores at or above the hand-code arm, or the contrast is non-significant; the pre-registered interpretation tree then reports all intervention contrasts as estimation-only, and the replication failure is itself a headline finding given the base trial's prominence.

\textbf{Hypothesis (plain language).} The published skill-formation gap is real enough to reappear in an independent replication of its own two arms.

\textbf{Metrics.} Quiz score by arm; standardized gap versus the published effect size.

\textbf{Result.} \emph{[placeholder --- fill in after experiments are run]}

\subsection{RQ8: Does Nudge Content Help Beyond Reactivity, Without Collateral Cost?}

\textbf{Setup.} Primary contrast: active-nudge vs.\ placebo arms, intention-to-treat and within the pre-trigger flagged subgroup (risk frozen before notification, so the subgroup is pre-treatment in all AI arms); secondary contrasts: placebo vs.\ unnotified (reactivity) and active vs.\ unnotified (total policy effect); guardrails (completion, time-on-task within 10\%) against the unnotified arm; 100\% intervention-integrity audit of firings; detection lag logged in all AI arms. The pilot is pre-registered as estimation-plus-gate: report the confidence interval on the content contrast and proceed to scale-up ($n \approx 120$--160, micro-randomized timing per \citealp{klasnja2015mrt}) iff the point estimate is at least 40\% of the internally replicated gap with guardrails clean. The 40\%-of-gap floor is an authored target: recovering less than ${\sim}7$ of 17 points would not justify a deployed interruption channel under the annoyance budget.

\textbf{Expected outcome.} Active exceeds placebo with a point estimate of at least 40\% of the replicated gap; guardrails hold; the reactivity contrast is small.

\textbf{Falsification condition.} Refuted if active does not exceed placebo (content adds nothing beyond reactivity), or if guardrails fail (the nudge ``works'' by degrading the session --- vacuous), or if the effect appears only outside the flagged subgroup (misattributed to the trigger); each branch's reporting language is pre-committed in the interpretation tree.

\textbf{Hypothesis (plain language).} What the nudge says --- not the fact that it interrupts --- buys back part of the lost comprehension, without slowing the work.

\textbf{Metrics.} Quiz score by arm (ITT and flagged subgroup); completion rate; time-on-task; integrity rate; detection lag.

\textbf{Result.} \emph{[placeholder --- fill in after experiments are run]}

\paragraph{Further evaluation.}
With more resources, natural extensions are pattern-matched nudge variants (a different question for a debugging-loop session than for a full-delegation one), a second task/library to test generalization beyond Trio, and micro-randomized timing designs \citep{klasnja2015mrt} for the intervention-timing question --- all deferred to scale-up, all inheriting the frozen machinery.

\subsection{Threats to Validity}

\paragraph{Construct validity.}
The erosive/protective construct is inherited from the published protocol; its reliability at this site is gated ($\kappa \geq 0.7$ on the binary before any model training), and the six-way taxonomy is demoted to descriptive use. The LLM intent tagger is an instrument inside the instrument: it runs under a frozen prompt with a 15\% human audit at $\kappa \geq 0.75$, and if it fails the audit, the deterministic-only feature set is the declared instrument (RQ4 measures the cost of that fallback). The prefix ceiling itself depends on annotators not inferring endings from prefixes; blinding to session endings and the prefix-time validity check (flagged sessions must differ in subsequent behavior) address this.

\paragraph{Internal validity.}
The intervention contrasts are protected by three mechanisms: matched placebo (content vs.\ reactivity), pre-notification risk freezing (flagged subgroups defined pre-treatment), and additive-only nudges with completion/time guardrails (no degraded-assistant confound). The offline stage's leakage controls are structural: prefix features cannot see the future by construction, cross-fitting is developer-disjoint, the cutoff is chosen on fold-train data only, and the gate set is sealed and scored once by a hashed, frozen model.

\paragraph{External validity.}
Findings are scoped to one task family (unfamiliar-library learning, the Trio setting or its documented reconstruction), one population (crowd-recruited developers with 1+ year Python, no Trio), and one intervention class (one-shot additive comprehension nudges). The base paper's own caveat that its setup differs from agentic coding products is inherited: results are a plausible lower bound for settings with richer delegation affordances, and no claim is made beyond the tested grid. Crowd-population drift from the original trial's population is bounded by screening traps and documented deviations, not eliminated.

\paragraph{Conclusion validity.}
All offline inference is cluster-bootstrap over developers with the full pipeline re-run per replicate; the gate uses exact binomial bounds \citep{clopper1934use} on sealed data; the live pilot is confirmatorily powered only for the replication contrast and is pre-registered as estimation-plus-gate for the intervention contrast, with the 40\%-of-gap decision rule and the interpretation tree fixed in advance. Attrition and missingness are handled by 20\% over-recruitment, multiple imputation \citep{rubin1987multiple}, and MNAR tipping-point sensitivity rather than complete-case optimism.

\section{Related Work}

The closest prior work is \citet{shen2026skill}'s skill-formation RCT --- whose task, quiz rubric, labels, and two arms this design reuses and extends --- and, for the detection claim specifically, \citet{yao2026scaffold}'s ``AI-prompt signature''; it is worth stating the delta plainly before surveying the landscape: the former labels complete sessions in hindsight and never estimates anything at prefix time, the latter detects a signature retrospectively at cohort level without an individual-session model or any intervention. The closed detect$\to$intervene$\to$verify loop on a learning outcome does have a precedent --- purpose-built intelligent tutors have run it for two decades, most sharply in the gaming-the-system line conceded below --- but no work closes that loop where the sensed stream is a general-purpose assistant's query log, the setting is authentic work, and the outcome is skill formation under the assistant causing the risk.

\paragraph{The causal premise and its neighbors.}
\citet{shen2026skill} provide the base RCT: 52 developers, a 17-point quiz gap without time savings, and six hindsight-coded interaction patterns that track the outcome. Their own future work looks at longitudinal dissipation, not within-session detection --- the timing question is stated by their framing but not attempted. \citet{kosmyna2025brain} corroborate the direction in essay writing with EEG, recall, and ownership measures --- different domain, different instrument, a preprint cited as convergent direction only. \citet{becker2025measuring} found experienced open-source maintainers \emph{slowed} by early-2025 AI tools --- a productivity result on experts that undercuts ``AI assistance is an unqualified win'' but says nothing about learning or prefixes. \citet{liu2026persistence} add a convergent causal datapoint from a different task family: AI assistance reduces persistence and hurts subsequent independent performance --- same direction as the base trial, again with no session-level detection and no intervention. \citet{aouad2026paradox} supply the theoretical ``why early matters'': heterogeneous AI literacy compounds into skill polarization, and if the compounding runs through early interaction differences, early intervention is the lever their model predicts --- their model, our motivating inference.

\paragraph{Behavioral signatures of AI reliance.}
\citet{yao2026scaffold} detects an AI-prompt signature (high first-attempt acceptance, few debugging retries) in competitive-programming platform data and shows that AI-prohibited evaluation gates screen for durable skill --- the strongest existence proof that the signal lives in behavior, but observational, cohort-level, retrospective, and open-loop. \citet{cao2022reliance} find a real-time behavioral proxy (gaze duration) more informative than retrospective self-report for reliance in AI-assisted decision-making --- a different sensor (hardware eye tracking) and a different outcome (decision accuracy); whether zero-hardware query-log features carry comparable signal for skill-formation risk is untested by anyone, and is exactly what RQ4 measures. (Engagement disclosure: this work is search-confirmed but not full-text-read this cycle --- ACM DL access failed --- and carries a pre-submission reading debt.) The detector's \emph{shape} --- a terminal learning outcome predicted from a partial behavioral log stream with counting-based features at early checkpoints --- is itself inherited from the learning-analytics early-warning tradition, exemplified by MOOC dropout and success prediction \citep{gardner2018mooc}; we concede the shape and claim what that literature does not occupy: the estimand (an erosive AI-interaction pattern, not course attrition), the decision clock (queries and minutes within one session, not weeks in a course), and the measured human prefix-labeling ceiling, for which early-warning work has no analogue.

\paragraph{Just-in-time adaptive interventions: the loop template, credited.}
The mobile-health JITAI literature formalized decision points, tailoring variables, and proximal versus distal outcomes \citep{nahumshani2018jitai}, and micro-randomized trials for intervention timing \citep{klasnja2015mrt}. No JITAI work senses AI-interaction patterns or targets skill formation; we import the vocabulary --- tailoring variable: the running risk score; proximal outcome: post-nudge pattern shift; distal outcome: quiz score --- rather than reinventing it, and the micro-randomized design is the pre-committed scale-up vehicle for timing questions. Our novelty claim is scoped accordingly: outside purpose-built tutoring systems --- whose closed-loop precedent is conceded in the next paragraph --- the first closed JITAI-style loop whose sensed state is an AI-interaction pattern during authentic work and whose distal outcome is measured skill formation.

\paragraph{Appropriate reliance and cognitive forcing.}
\citet{bucinca2021trust} showed that cognitive forcing functions --- interruption-based nudges --- reduce overreliance on AI in decision-making, at a user-experience cost; this is the nearest intervention class with RCT evidence, and their finding that effects partly run through friction is a direct design input here: it is why the placebo arm exists and why the gate carries an explicit annoyance budget rather than treating false positives as free. The automation-bias tradition \citep{lee2004trust} frames why delegation is the assistant's default happy path --- the background condition that makes erosive patterns the path of least resistance.

\paragraph{Learner modeling and the design-principles response.}
Intelligent tutoring systems have run closed-loop learner modeling for three decades --- Bayesian knowledge tracing \citep{corbett1995knowledge} and its successors --- and the sharpest prior twin of our loop lives here: the gaming-the-system line (Baker and colleagues, 2004--2008), whose ``Scooter the Tutor'' deployment detected a maladaptive \emph{interaction pattern} --- not curricular mastery --- from log features, intervened live on detected students while sparing the rest, and measured learning outcomes: a closed detect$\to$intervene$\to$verify loop on a learning outcome, two decades before this design. We concede that precedent by name rather than claim the loop. What it does not occupy is this setting: its sensed state is exploitation of a purpose-built tutor's own help channel, whereas here the ``tutor'' is a general-purpose assistant that is itself the erosion vector, the work is authentic rather than curricular, and the authorization machinery --- a measured human prefix ceiling and a sealed prevalence-anchored gate --- has no analogue there, because a tutor intervening on its own students never needed a licence to interrupt. The nearest real-time education RCT is Tutor CoPilot \citep{wang2024tutorcopilot}, which delivers expert-like guidance into live tutoring sessions --- but its sensed state is tutoring context, not a learner's AI-interaction pattern, and its outcome is tutoring quality, not skill formation under the assistant itself. Finally, Agents-That-Teach \citep{mehra2026agents} represents the design-principles answer to the same harm: build incidental learning back into AI-assisted development via pedagogical scaffolding, prescribed open-loop --- applied to all users, unconditioned on any measured learner state. That stance taxes users already in protective patterns (the 65--86\% group needs no Socratic gate) and cannot tell whether its principles are working mid-session. We are the closed-loop complement: measure first, intervene conditionally, verify on the same quiz that established the harm --- and if their designs prove adaptive on full reading, the pre-committed repositioning is that we supply the measurement layer their adaptivity would require.

\section{Conclusion}

Skill erosion under AI assistance is causally established, pattern-mediated, and --- so far --- visible only after the quiz. This paper's wager is that the drift is catchable in the query stream while the session is still running, and it insists that the claim is only worth making under discipline: a prognostic target instead of an ill-posed prefix-time diagnosis, a measured human ceiling instead of an assumed one, an exact, sealed, prevalence-anchored gate before any human is intervened on, a matched placebo that eats the reactivity story, and an internal replication that keeps the borrowed baseline honest. The design is serialized so that evidence, not optimism, decides how far it goes --- and either branch produces knowledge the field lacks: how early erosive patterns are visible at all, or whether a one-shot additive nudge can buy back part of a 17-point comprehension loss without taxing the users who never needed it. Beyond this study, the instrument itself --- prefix-only, model-internal-free risk scoring with an auditable authorization gate --- is a template for measuring, at decision time, harms that today are only measurable after they are locked in. The measured prefix ceiling and detection-lag distribution will also bound what any future real-time protection can promise, a constraint the design-principles literature currently has no way to state.

\bibliographystyle{plainnat}
\bibliography{refs}

\end{document}
